% ============================================================
% INTRODUCTION
% ============================================================
\chapter*{Introduction}
\pagestyle{styleGlobal}
\label{sec:introduction}
\markboth{Introduction}{Introduction}
\addcontentsline{toc}{chapter}{Introduction}

Le numérique a changé la façon dont les personnes accèdent aux services professionnels. Aujourd'hui, des plateformes numériques permettent de trouver facilement un professionnel compétent pour réaliser une tâche précise, de consulter son profil, de le contacter et de le payer en toute sécurité. Cependant, dans bien des environnements, cette mise en relation entre prestataires et clients reste informelle, peu structurée et difficile à organiser. C'est dans ce contexte qu'Akieni, société spécialisée dans la transformation digitale au Congo, a exprimé le besoin de disposer d'une solution numérique capable de connecter de manière fiable des prestataires de services professionnels avec des clients. C'est ainsi qu'est né «~SALISA~», notre solution pour répondre directement à ce besoin.

La problématique principale de ce travail peut être formulée ainsi : comment concevoir et réaliser une application cross-plateforme qui permet de connecter efficacement des prestataires et des clients, en garantissant la qualité des services, la sécurité des transactions et une expérience simple pour tous les utilisateurs ?

Pour répondre à cette problématique, nous formulons les hypothèses suivantes. Premièrement, un système de vérification d'identité et de notation permettra d'établir la confiance entre les utilisateurs de la plateforme. Deuxièmement, un algorithme de mise en correspondance basé sur les compétences, la localisation et le budget aidera les clients à trouver rapidement le bon prestataire. Troisièmement, l'intégration du paiement Mobile Money avec un mécanisme de sécurisation des fonds (escrow) protégera les deux parties lors des transactions. Quatrièmement, l'utilisation du processus 2TUP combiné au langage de modélisation UML permettra de concevoir un système solide et bien structuré.

Ce travail s'intitule : \textbf{«~Conception et réalisation d'une application cross-plateforme de mise en relation entre prestataires et clients pour services professionnels à Akieni~»}. Il est organisé en trois grandes parties. La première présente la structure d'accueil Akieni ainsi que le cadrage général du travail. La deuxième est consacrée à l'analyse des besoins et à la conception du système selon la démarche 2TUP et UML. La troisième décrit les outils utilisés, le planning et les résultats obtenus.
